% engine=luatex language=uk

% lua.newtable

\chapter{与 \TEX 相关的库}\label{tex}

\section{\type {lua}}

\subsection{版本信息}

\topicindex{libraries+\type{lua}}
\topicindex{version}

\libindex{version}

此库包含一个只读项目：

\begin{typing}
<string> s = lua.version
\end{typing}

这会返回 \LUA 版本标识字符串。当前值为
\directlua {tex.print(lua.version)}。

\section{\type {status}}

\topicindex{libraries+\type{status}}

\libindex{list}
\libindex{resetmessages}
\libindex{setexitcode}

此库包含一些运行时配置项，你可能会发现它们在消息报告中很有用；以及一个
迭代器函数，它以 table 形式获取所有名称和值。

\begin{typing}
<table> info = status.list()
\end{typing}

table 中的键是已知项目，值为当前值。\type {status} 中几乎所有值都是在运行时
通过元表在访问时获取的，因此你不能对 \type {status} 使用 \type {pairs}，但
你{\it 可以}对 \type {info} 使用 \type {pairs}。如果你不需要完整列表，也可以
通过将名称作为 \type {status} 的索引来请求单个项目。当前列表如下：

\begin{longtable}{l p{5cm}}
\toprule 键                       & 说明 \\
\midrule
\type{banner}             & 终端显示横幅 \\
\type{best_page_break}    & 当前最佳断点（一个节点） \\
\type{buf_size}           & 当前分配的行缓冲区大小 \\
\type{callbacks}          & 迄今为止执行的回调总数 \\
\type{cs_count}           & 控制序列数量 \\
\type{dest_names_size}    & \PDF 目标表大小 \\
\type{dvi_gone}           & 已写入的 \DVI 字节数 \\
\type{dvi_ptr}            & 尚未写入的 \DVI 字节数 \\
\type{dyn_used}           & 正在使用的令牌（多字）内存  \\
\type{filename}           & 当前输入文件名 \\
\type{fix_mem_end}        & 已使用令牌的最大数量 \\
\type{fix_mem_min}        & 为令牌分配的最小字数 \\
\type{fix_mem_max}        & 为令牌分配的最大字数 \\
\type{font_ptr}           & 活动字体数量 \\
\type{hash_extra}         & 额外允许的哈希量 \\
\type{hash_size}          & 哈希表大小 \\
\type{indirect_callbacks} & 其中本身是其他回调（例如文件读取器）结果的数量 \\
\type{ini_version}        & 如果这是一次 initex 运行则为 \type {true} \\
\type{init_pool_ptr}      & initex 字符串池索引 \\
\type{init_str_ptr}       & initex 字符串数量 \\
\type{input_ptr}          & 我们当前所在的输入层级 \\
\type{inputid}            & 当前输入的数字 id \\
\type{largest_used_mark}  & 最大引用的 marks 类 \\
\type{lasterrorcontext}   & 最后一个错误上下文字符串（含换行） \\
\type{lasterrorstring}    & 最后一个 \TEX 错误字符串 \\
\type{lastluaerrorstring} & 最后一个 \LUA 错误字符串 \\
\type{lastwarningstring}  & 最后一个警告标签，通常指示位于哪个部分\\
\type{lastwarningtag}     & 最后一个警告字符串\\
\type{linenumber}         & 当前输入文件中的位置 \\
\type{log_name}           & 日志文件名 \\
\type{luabytecode_bytes}   & \LUA 字节码寄存器中的字节数 \\
\type{luabytecodes}        & 活动 \LUA 字节码寄存器数量 \\
\type{luastate_bytes}      & \LUA 解释器正在使用的字节数 \\
\type{luatex_engine}       & \LUATEX 引擎标识符 \\
\type{luatex_hashchars}    & \LUA 对字符串做哈希的长度（$2^n$） \\
\type{luatex_hashtype}    & 使用的哈希方法（在 \LUAJITTEX 中） \\
\type{luatex_version}      & \LUATEX 版本号 \\
\type{luatex_revision}     & \LUATEX 修订字符串 \\
\type{max_buf_stack}       & 最大已用缓冲区位置 \\
\type{max_in_stack}        & 最大已用输入栈条目 \\
\type{max_nest_stack}      & 最大已用嵌套栈条目 \\
\type{max_param_stack}     & 最大已用参数栈条目 \\
\type{max_save_stack}      & 最大已用保存栈条目 \\
\type{max_strings}         & 最大允许字符串数 \\
\type{nest_size}           & 嵌套栈大小 \\
\type{node_mem_usage}      & 一个字符串，提供当前使用节点的洞察 \\
\type{obj_ptr}             & 最大 \PDF 对象指针 \\
\type{obj_tab_size}        & \PDF 对象表大小 \\
\type{output_active}       & 如果 \prm {output} 例程处于活动状态则为 \type {true} \\
\type{output_file_name}    & \PDF 或 \DVI 文件名 \\
\type{param_size}          & 参数栈大小 \\
\type{pdf_dest_names_ptr}  & 最大 \PDF 目标指针 \\
\type{pdf_gone}            & 已写入的 \PDF 字节数 \\
\type{pdf_mem_ptr}         & 最大已用 \PDF 内存 \\
\type{pdf_mem_size}        & \PDF 内存大小 \\
\type{pdf_os_cntr}         & 最大 \PDF 对象流指针 \\
\type{pdf_os_objidx}       & \PDF 对象流索引 \\
\type{pdf_ptr}             & 尚未写入的 \PDF 字节数 \\
\type{pool_ptr}           & 字符串池索引 \\
\type{pool_size}           & 当前为字符串字符分配的大小 \\
\type{save_size}           & 保存栈大小 \\
\type{shell_escape}        & \type {0} 表示禁用，\type {1} 表示允许任何操作，\type {2} 为受限 \\
\type{safer_option}        & \type {1} 表示强制执行更安全模式 \\
\type{luadebug_option}     & 如果在命令行启用了 \type{debug} 库则为 true \\
\type{output_directory}    & 输出目录的值 \\
\type{kpse_used}           & \type {1} 表示使用 kpse \\
\type{stack_size}          & 输入栈大小 \\
\type{str_ptr}             & 字符串数量 \\
\type{total_pages}         & 已写入的页数 \\
\type{var_mem_max}         & 为节点分配的字数 \\
\type{var_used}            & 正在使用的变量（单字）内存 \\
\type{lc_collate}          & 启动时 \type {LC_COLLATE} 的值（启动时变为 \type {C}） \\
\type{lc_ctype}            & 启动时 \type {LC_CTYPE} 的值（启动时变为 \type {C}） \\
%NC \type{lc_monetary}         & 启动时 \type {LC_MONETARY} 的值 \\
\type{lc_numeric}           & 启动时 \type {LC_NUMERIC} 的值 \\
%NC \type{lc_time}             & 启动时 \type {LC_TIME} 的值（启动时变为 \type {C}） \\
\bottomrule
\end{longtable}

错误和警告消息可以通过 \type {resetmessages} 函数清除。可以使用
\type {setexitcode} 设置返回值。

\section{\type {tex}}

\subsection{简介}

\topicindex{libraries+\type{tex}}

\type {tex} table 包含大量虚拟内部 \TEX 参数，其中部分可写。

``虚拟这一称谓意味着这些项目并非在 \LUA 中正式定义，而只是由元表
处理实际 \TEX 值的前端。因此，大多数 \LUA table 操作符（如 \type {pairs} 和
\type {#}）对此类项目不起作用。

目前，几乎可以访问所有你能在 \prm {the} 之后使用的参数，即单个令牌或在
\TEX 中具有特殊性质的参数。这不包括需要额外参数的参数，如 \type {\the\scriptfont}。
由简单整数和尺寸寄存器组成的子集既可读又可写（如 \prm {tracingcommands}
和 \prm {parindent}）。

\subsection{内部参数值：\type {set}}

\topicindex{parameters+internal}

\libindex{set}
\libindex{get}

对于本节中的所有参数，可以直接使用其名称作为 \type {tex} table 中的索引来
访问，或者使用 \type {tex.get} 和 \type {tex.set} 函数之一。

确切的参数和返回值因实际参数而异，\type {tex.set} 是否有作用也是如此。对于
{\em 可以}设置的参数，可以使用 \type {global} 作为 \type {tex.set} 的第一个
参数；这使赋值变为全局而非局部。

\begin{typing}
tex.set (["global",] <string> n, ...)
... = tex.get (<string> n)
\end{typing}

黏胶有点特殊，因为涉及五个值。返回值是一个 \nod {glue_spec} 节点，但当你向
\type {tex.get} 传递 \type {false} 作为最后一个参数时，你得到黏胶的宽度；
传递 \type {true} 则得到全部五个值。否则你得到一个节点，它是内部值的副本，
因此你有责任在 \LUA 端释放它。设置黏胶量时，你可以传递一个 \nod {glue_spec}
或最多五个数字。向 \type {get} 传递 \type {true} 时，黏胶返回 5 个值；
传递 \type {false} 时仅返回宽度。

\subsubsection{整数参数}

整数参数接受并返回 \LUA 数字。以下为可读写的：

\starttwocolumns
\begin{typing}
tex.adjdemerits
tex.binoppenalty
tex.brokenpenalty
tex.catcodetable
tex.clubpenalty
tex.day
tex.defaulthyphenchar
tex.defaultskewchar
tex.delimiterfactor
tex.displaywidowpenalty
tex.doublehyphendemerits
tex.endlinechar
tex.errorcontextlines
tex.escapechar
tex.exhyphenpenalty
tex.fam
tex.finalhyphendemerits
tex.floatingpenalty
tex.globaldefs
tex.hangafter
tex.hbadness
tex.holdinginserts
tex.hyphenpenalty
tex.interlinepenalty
tex.language
tex.lastlinefit
tex.lefthyphenmin
tex.linepenalty
tex.localbrokenpenalty
tex.localinterlinepenalty
tex.looseness
tex.mag
tex.maxdeadcycles
tex.month
tex.newlinechar
tex.outputpenalty
tex.pausing
tex.postdisplaypenalty
tex.predisplaydirection
tex.predisplaypenalty
tex.pretolerance
tex.relpenalty
tex.righthyphenmin
tex.savinghyphcodes
tex.savingvdiscards
tex.showboxbreadth
tex.showboxdepth
tex.time
tex.tolerance
tex.tracingassigns
tex.tracingcommands
tex.tracinggroups
tex.tracingifs
tex.tracinglostchars
tex.tracingmacros
tex.tracingnesting
tex.tracingonline
tex.tracingoutput
tex.tracingpages
tex.tracingparagraphs
tex.tracingrestores
tex.tracingscantokens
tex.tracingstats
tex.uchyph
tex.vbadness
tex.widowpenalty
tex.year
\end{typing}
\stoptwocolumns

以下为只读的：

\startthreecolumns
\begin{typing}
tex.deadcycles
tex.insertpenalties
tex.parshape
tex.interlinepenalties
tex.clubpenalties
tex.widowpenalties
tex.displaywidowpenalties
tex.prevgraf
tex.spacefactor
\end{typing}
\stopthreecolumns

\subsubsection{尺寸参数}

尺寸参数接受 \LUA 数字（表示缩放点）或字符串（含尺寸单位）。结果始终为以
缩放点为单位的数字。以下为可读写的：

\startthreecolumns
\begin{typing}
tex.boxmaxdepth
tex.delimitershortfall
tex.displayindent
tex.displaywidth
tex.emergencystretch
tex.hangindent
tex.hfuzz
tex.hoffset
tex.hsize
tex.lineskiplimit
tex.mathsurround
tex.maxdepth
tex.nulldelimiterspace
tex.overfullrule
tex.pagebottomoffset
tex.pageheight
tex.pageleftoffset
tex.pagerightoffset
tex.pagetopoffset
tex.pagewidth
tex.parindent
tex.predisplaysize
tex.scriptspace
tex.splitmaxdepth
tex.vfuzz
tex.voffset
tex.vsize
tex.prevdepth
tex.prevgraf
tex.spacefactor
\end{typing}
\stopthreecolumns

以下为只读的：

\startthreecolumns
\begin{typing}
tex.pagedepth
tex.pagefilllstretch
tex.pagefillstretch
tex.pagefilstretch
tex.pagegoal
tex.pageshrink
tex.pagestretch
tex.pagetotal
\end{typing}
\stopthreecolumns

注意：与所有 \LUA table 一样，你可以向其中添加值。因此，以下是有效的：

\begin{typing}
tex.foo = 123
\end{typing}

当你访问一个 \TEX 参数时，会进行查找。对于只读变量，这意味着你会得到返回值，
但当你设置它们时，会在 table 中创建一个新条目，从而使原始条目不可见。

有几个特殊情况我们做了例外处理：\type {prevdepth}、\type {prevgraf} 和
\type {spacefactor}。通常它们通过 \type {tex.nest} table 访问：

\begin{typing}
tex.nest[tex.nest.ptr].prevdepth   = p
tex.nest[tex.nest.ptr].spacefactor = s
\end{typing}

不过，以下方式也有效：

\begin{typing}
tex.prevdepth   = p
tex.spacefactor = s
\end{typing}

请记住，当你在 \LUA 端直接操作节点列表时，可能需要显式更新嵌套栈顶的
\type {prevdepth}，因为 \LUATEX 无法猜测你的意图。通过使用 \type {tex}
table 中的访问器，你可以获取和设置嵌套栈顶的值。

\subsubsection{方向参数}

方向参数为只读，返回一个 \LUA 字符串。

\startthreecolumns
\begin{typing}
tex.bodydir
tex.mathdir
tex.pagedir
tex.pardir
tex.textdir
\end{typing}
\stopthreecolumns

\subsubsection{黏胶参数}

黏胶参数接受并返回一个表示 \nod {glue_spec} 节点的 userdata 对象。

\startthreecolumns
\begin{typing}
tex.abovedisplayshortskip
tex.abovedisplayskip
tex.baselineskip
tex.belowdisplayshortskip
tex.belowdisplayskip
tex.leftskip
tex.lineskip
tex.parfillskip
tex.parskip
tex.rightskip
tex.spaceskip
tex.splittopskip
tex.tabskip
tex.topskip
tex.xspaceskip
\end{typing}
\stopthreecolumns

\subsubsection{muglue 参数}

所有 muglue 参数均为只读，返回一个 \LUA 字符串。

\startthreecolumns
\begin{typing}
tex.medmuskip
tex.thickmuskip
tex.thinmuskip
\end{typing}
\stopthreecolumns

\subsubsection{令牌列表参数}

令牌列表参数接受并返回 \LUA 字符串。\LUA 字符串使用 \prm {the} \prm {toks}
式展开与令牌列表相互转换：所有类别码要么是空格（10），要么是其他（12）。
因此，对其中一些参数（如 ``tex.output）赋值实际上没有用处，但考虑到
未来将支持完整令牌字符串的扩展，做例外处理并不合适。

\startthreecolumns
\begin{typing}
tex.errhelp
tex.everycr
tex.everydisplay
tex.everyeof
tex.everyhbox
tex.everyjob
tex.everymath
tex.everypar
tex.everyvbox
tex.output
\end{typing}
\stopthreecolumns

\subsection{转换命令}

\topicindex{convert commands}

所有 ``convert 命令均为只读，返回一个 \LUA 字符串。当前支持的命令有：

\starttwocolumns
\begin{typing}
tex.eTeXVersion
tex.eTeXrevision
tex.formatname
tex.jobname
tex.luatexbanner
tex.luatexrevision
tex.fontname(number)
tex.uniformdeviate(number)
tex.number(number)
tex.romannumeral(number)
tex.fontidentifier(number)
\end{typing}
\stoptwocolumns

如果你想知道为什么这个列表看起来杂乱无章；这些都是不需要参数的 ``convert 内部命令的所有情况，以及只需要简单数值的命令。特殊的（仅限 \LUA）
\type {tex.fontidentifier} 情况返回与字体 id 编号匹配的 \type {csname} 字符串
（如果存在的话）。

\subsection{最后一项命令}

\topicindex{last items}

所有 ``last item 命令均为只读，返回一个数字。当前支持的命令有：

\startthreecolumns
\begin{typing}
tex.lastpenalty
tex.lastkern
tex.lastskip
tex.lastnodetype
tex.inputlineno
tex.lastxpos
tex.lastypos
tex.randomseed
tex.luatexversion
tex.eTeXminorversion
tex.eTeXversion
tex.currentgrouplevel
tex.currentgrouptype
tex.currentiflevel
tex.currentiftype
tex.currentifbranch
\end{typing}
\stopthreecolumns

\subsection{访问寄存器：\type {set*}}

\topicindex{attributes}
\topicindex{registers}

\libindex{attribute}  \libindex{setattribute}  \libindex{getattribute}  \libindex{isattribute}
\libindex{count}      \libindex{setcount}      \libindex{getcount}      \libindex{iscount}
\libindex{dimen}      \libindex{setdimen}      \libindex{getdimen}      \libindex{isdimen}
\libindex{skip}       \libindex{setskip}       \libindex{getskip}       \libindex{isskip}
\libindex{muskip}     \libindex{setmuskip}     \libindex{getmuskip}     \libindex{ismuskip}
\libindex{glue}       \libindex{setglue}       \libindex{getglue}       \libindex{isglue}
\libindex{muglue}     \libindex{setmuglue}     \libindex{getmuglue}     \libindex{ismuglue}
\libindex{toks}       \libindex{settoks}       \libindex{gettoks}       \libindex{istoks}
\libindex{box}        \libindex{setbox}        \libindex{getbox}        \libindex{isbox}

\libindex{scantoks}

\libindex{getmark}

\TEX 的属性（\lpr {attribute}）、计数器（\prm {count}）、尺寸（\prm {dimen}）、
skip（\prm {skip}、\prm {muskip}）和令牌（\prm {toks}）寄存器可以通过
\type {tex} table 的两组各五个虚拟子表来访问和写入：

\startthreecolumns
\begin{typing}
tex.attribute
tex.count
tex.dimen
tex.skip
tex.glue
tex.muskip
tex.muglue
tex.toks
\end{typing}
\stopthreecolumns

可以将相关的 \lpr {attributedef}、\prm {countdef}、\prm {dimendef}、
\prm {skipdef} 或 \prm {toksdef} 控制序列的名称用作这些 table 的索引：

\begin{typing}
tex.count.scratchcounter = 0
enormous = tex.dimen['maxdimen']
\end{typing}

在这种情况下，\LUATEX 会即时为你查找值。你必须使用有效的 \prm {countdef}
（或 \lpr {attributedef}、\prm {dimendef}、\prm {skipdef}、\prm {toksdef}），
其他任何内容都会产生错误（其意图是最终也允许 \type {<chardef tokens>}
甚至展开为数字的宏）。

\begin{itemize}\item count 寄存器接受并返回 \LUA 数字。
    

    \item 尺寸寄存器接受 \LUA 数字（以缩放点为单位）或字符串（含绝对尺寸；
        \type {em}、\type {ex} 和 \type {px} 不允许）。结果始终为以缩放点
        为单位的数字。
    

    \item 令牌寄存器接受并返回 \LUA 字符串。\LUA 字符串使用 \prm {the}
        \prm {toks} 式展开与令牌列表相互转换：所有类别码要么是空格（10），
        要么是其他（12）。
    

    \item skip 寄存器接受并返回 \nod {glue_spec} userdata 节点对象（参见本手册
        其他部分对节点接口的描述）。
    

    \item glue 寄存器就是 skip 寄存器，但不是 userdata 形式而是详细表示形式。
    

    \item 与 count 一样，属性寄存器接受并返回 \LUA 数字。
    

\end{itemize}

除了数组寻址之外，还为所有情况定义了访问器函数，例如，以下是 \prm {skip}
寄存器的一组可能函数：

\begin{typing}
tex.setskip (["global",] <number> n, <node> s)
tex.setskip (["global",] <string> s, <node> s)
<node> s = tex.getskip (<number> n)
<node> s = tex.getskip (<string> s)
\end{typing}

我们对 \type {count}、\type {dimen}、\type {muskip} 和 \type {toks} 有类似
的设置函数。计数器和尺寸以数字表示，skip 和 muskip 以节点表示，toks 以字符串
表示。

同样，glue 变体不使用 \nod {glue_spec} userdata 节点。\type {setglue}
函数最多接受 5 个参数：宽度、伸展、收缩、伸展阶和收缩阶。如果你不传递值或
某个值不是数字，相应属性将变为零。\type {getglue} 函数报告所有属性，除非
第二个参数为 \type {false}，此时仅返回宽度。

以下是一个使用三元组的示例：

\begin{typing}
local d = tex.getdimen("foo")
if tex.isdimen("bar") then
    tex.setdimen("bar",d)
end
\end{typing}

有六个额外的 skip（glue）相关辅助函数：

\begin{typing}
tex.setglue (["global"], <number> n,
    width, stretch, shrink, stretch_order, shrink_order)
tex.setglue (["global"], <string> s,
    width, stretch, shrink, stretch_order, shrink_order)
width, stretch, shrink, stretch_order, shrink_order =
    tex.getglue (<number> n)
width, stretch, shrink, stretch_order, shrink_order =
    tex.getglue (<string> s)
\end{typing}

另外两个是 \type {tex.setmuglue} 和 \type {tex.getmuglue}。

\type {dimen}、\type {count}、\type {skip}、\type {muskip}、\type {box}
和 \type {attribute} 寄存器都有这样的辅助函数，但 glue 的比较特殊，因为
它们需要处理更多属性。

与之前讨论的通用 \type {get} 和 \type {set} 函数一样，对于 skip 寄存器，
\type {getskip} 返回节点，\type {getglue} 返回数字；而 \type {setskip}
接受节点，\type {setglue} 期望最多五个数字。同样，当你向 \type {getglue}
传递 \type {false} 作为第二个参数时，仅返回宽度。\type {getmuskip}、
\type {setmuskip} 的 mu 变体也是如此。

对于令牌寄存器，我们有一个备选方案，可以指定类别码表：

\begin{typing}
tex.scantoks(0,3,"$e=mc^2$")
tex.scantoks("global",0,"$\int\limits^1_2$")
\end{typing}

在基于函数的接口中，可以通过使用字符串 \type {global} 作为第一个函数参数
来全局定义值。

marks 有一个专用的获取函数：\type {getmark}，它接受两个参数。第一个参数是
\type {top}、\type {bottom}、\type {first}、\type {splitbottom} 或
\type {splitfirst} 之一，第二个参数是 marks 类编号。当不提供参数时，返回
当前最大类数。

\subsection{{字符码寄存器：\type {[get|set}*code[s]}}]

\topicindex{characters+codes}

\libindex{lccode}    \libindex{setlccode}    \libindex{getlccode}
\libindex{uccode}    \libindex{setuccode}    \libindex{getuccode}
\libindex{sfcode}    \libindex{setsfcode}    \libindex{getsfcode}
\libindex{catcode}   \libindex{setcatcode}   \libindex{getcatcode}
\libindex{mathcode}  \libindex{setmathcode}  \libindex{getmathcode}
\libindex{delcode}   \libindex{setdelcode}   \libindex{getdelcode}

\libindex{setdelcodes}   \libindex{getdelcodes}
\libindex{setmathcodes}  \libindex{getmathcodes}

\TEX 的字符码表（\prm {lccode}、\prm {uccode}、\prm {sfcode}、\prm
{catcode}、\prm {mathcode}、\prm {delcode}）可以通过 \type {tex} table 的
六个虚拟子表来访问和写入

\startthreecolumns
\begin{typing}
tex.lccode
tex.uccode
tex.sfcode
tex.catcode
tex.mathcode
tex.delcode
\end{typing}
\stopthreecolumns

函数调用接口大致如上，但有一些曲折。\type {sfcode} 是最简单的：

\begin{typing}
tex.setsfcode (["global",] <number> n, <number> s)
<number> s = tex.getsfcode (<number> n)
\end{typing}

\type {lccode} 和 \type {uccode} 的函数调用接口还允许你同时设置关联的兄弟码：

\begin{typing}
tex.setlccode (["global"], <number> n, <number> lc)
tex.setlccode (["global"], <number> n, <number> lc, <number> uc)
<number> lc = tex.getlccode (<number> n)
tex.setuccode (["global"], <number> n, <number> uc)
tex.setuccode (["global"], <number> n, <number> uc, <number> lc)
<number> uc = tex.getuccode (<number> n)
\end{typing}

\type {catcode} 的函数调用接口还允许你在赋值或查询时指定要使用的类别表
（两种情况的默认值都是当前表）：

\begin{typing}
tex.setcatcode (["global"], <number> n, <number> c)
tex.setcatcode (["global"], <number> cattable, <number> n, <number> c)
<number> lc = tex.getcatcode (<number> n)
<number> lc = tex.getcatcode (<number> cattable, <number> n)
\end{typing}

\type {delcode} 和 \type {mathcode} 的接口使用小数组 table 来设置和检索值：

\begin{typing}
tex.setmathcode (["global"], <number> n, <table> mval )
<table> mval = tex.getmathcode (<number> n)
tex.setdelcode (["global"], <number> n, <table> dval )
<table> dval = tex.getdelcode (<number> n)
\end{typing}

其中 \type {mathcode} 的 table 是一个包含 3 个数字的数组，如下所示：

\begin{typing}
{
    <number> class,
    <number> family,
    <number> character
}
\end{typing}

而 \type {delcode} 的 table 是一个包含 4 个数字的数组，如下所示：

\begin{typing}
{
    <number> small_fam,
    <number> small_char,
    <number> large_fam,
    <number> large_char
}
\end{typing}

你也可以避免使用 table：

\begin{typing}
tex.setmathcode (["global"], <number> n, <number> class,
    <number> family, <number> character)
class, family, char =
    tex.getmathcodes (<number> n)
tex.setdelcode (["global"], <number> n, <number> smallfam,
    <number> smallchar, <number> largefam, <number> largechar)
smallfam, smallchar, largefam, largechar =
    tex.getdelcodes (<number> n)
\end{typing}

通常，根据 \lpr {Udelcode} 的用法，分隔符码赋值中的第三个和第四个值为零，
但返回的 table 中这些位置可能有值（例如，如果分隔符码是使用 \prm {delcode}
设置的）。未设置的 \type {delcode} 可以通过 \type {dval[1]} 为 $-1$ 来识别。

\subsection{{Box 寄存器：\type {[get|set}box}}]

\topicindex{registers}
\topicindex{boxes}

\libindex{box}
\libindex{setbox}  \libindex{getbox}

可以使用 \type {node} 库中定义的节点接口来设置和查询实际的 box，例如来自
\prm {hbox}、\prm {vbox} 或 \prm {vtop} 的 box：

\begin{typing}
tex.box
\end{typing}

用于数组访问，或

\begin{typing}
tex.setbox(["global",] <number> n, <node> s)
tex.setbox(["global",] <string> cs, <node> s)
<node> n = tex.getbox(<number> n)
<node> n = tex.getbox(<string> cs)
\end{typing}

用于基于函数的访问。在基于函数的接口中，可以通过使用字符串 \type {global}
作为第一个函数参数来全局定义值。

请注意，如下赋值

\begin{typing}
tex.box[0] = tex.box[2]
\end{typing}

不会复制节点列表，它只是复制了一个节点指针。如果 \type {\box2} 之后被
\TEX 命令清除，\type {\box0} 的内容也会变为无效。为防止这种情况发生，
始终使用 \type {node.copy_list}，除非你赋值给临时变量：

\begin{typing}
tex.box[0] = node.copy_list(tex.box[2])
\end{typing}

\subsection{{重用 box：\type {[use|save}boxresource} 与 \type {getboxresourcedimensions}}]

\topicindex{boxes+reuse}

\libindex{useboxresource}
\libindex{saveboxresource}
\libindex{getboxresourcedimensions}

以下函数将注册一个 box 以供重用（这是仿照 \PDF 中所谓的 xforms 建模的）。
你可以通过 \lpr {useboxresource} 或创建子类型为~2 的规则节点来（重）使用
该 box。

\begin{typing}
local index = tex.saveboxresource(n,attributes,resources,immediate,type,margin)
\end{typing}

可选的第二个和第三个参数是字符串，第四个是布尔值。第五个参数是类型。设置为
非零时省略 \type {/Type} 条目。值为 1 或 3 仍会写入 \type {/BBox}，而 2 或 3
会写入 \type {/Matrix}。第六个参数是（虚拟）边距，它超出 \TEX 所见的有效
包围盒。除了 box 编号之外，还可以传递一个 \type {[h|v]list} 节点。

你可以通过以下方式生成引用（规则类型）：

\begin{typing}
local reused = tex.useboxresource(n,wd,ht,dp)
\end{typing}

尺寸是可选的，最终尺寸作为额外值返回。以下只是一个附加功能（不返回尺寸意味着
该资源未知）：

\begin{typing}
local w, h, d, m = tex.getboxresourcedimensions(n)
\end{typing}

这会返回资源的宽度、高度、深度和边距。

\subsection{\type {triggerbuildpage}}

\topicindex{pages}

\libindex{triggerbuildpage}

你不应指望 \type {triggerbuildpage} 辅助函数有太多作为，因为通常当 \TEX
认为没有什么需要做时就不会做太多，但它可能对某些应用程序有用。它只是按照
其名称所做的那样——调用构建页面的内部函数，前提是有东西需要构建。

\subsection{\type {splitbox}}

\topicindex{boxes+split}

\libindex{splitbox}

你可以拆分一个 box：

\begin{typing}
local vlist = tex.splitbox(n,height,mode)
\end{typing}

余数保留在原始 box 中，并返回一个打包好的 vlist。此操作类似于 \prm {vsplit}
操作。mode 可以是 \type {additional} 或 \type {exactly}，涉及被拆分出去
的 box。

\subsection{{访问数学参数：\type {[get|set}math}}]

\topicindex{math+parameters}
\topicindex{parameters+math}

\libindex{setmath}
\libindex{getmath}

可以使用以下方式设置和查询内部数学参数：

\begin{typing}
tex.setmath(["global",] <string> n, <string> t, <number> n)
<number> n = tex.getmath(<string> n, <string> t)
\end{typing}

与之前一样，可选的第一个参数 \type {global} 表示全局赋值。

第一个字符串是去掉前导 ``Umath 的参数名，第二个字符串是去掉尾部
``style 的样式名。为完整起见，数学参数名的值如下：

\begin{typing}
quad                axis                operatorsize
overbarkern         overbarrule         overbarvgap
underbarkern        underbarrule        underbarvgap
radicalkern         radicalrule         radicalvgap
radicaldegreebefore radicaldegreeafter  radicaldegreeraise
stackvgap           stacknumup          stackdenomdown
fractionrule        fractionnumvgap     fractionnumup
fractiondenomvgap   fractiondenomdown   fractiondelsize
limitabovevgap      limitabovebgap      limitabovekern
limitbelowvgap      limitbelowbgap      limitbelowkern
underdelimitervgap  underdelimiterbgap
overdelimitervgap   overdelimiterbgap
subshiftdrop        supshiftdrop        subshiftdown
subsupshiftdown     subtopmax           supshiftup
supbottommin        supsubbottommax     subsupvgap
spaceafterscript    connectoroverlapmin
ordordspacing       ordopspacing        ordbinspacing     ordrelspacing
ordopenspacing      ordclosespacing     ordpunctspacing   ordinnerspacing
opordspacing        opopspacing         opbinspacing      oprelspacing
opopenspacing       opclosespacing      oppunctspacing    opinnerspacing
binordspacing       binopspacing        binbinspacing     binrelspacing
binopenspacing      binclosespacing     binpunctspacing   bininnerspacing
relordspacing       relopspacing        relbinspacing     relrelspacing
relopenspacing      relclosespacing     relpunctspacing   relinnerspacing
openordspacing      openopspacing       openbinspacing    openrelspacing
openopenspacing     openclosespacing    openpunctspacing  openinnerspacing
closeordspacing     closeopspacing      closebinspacing   closerelspacing
closeopenspacing    closeclosespacing   closepunctspacing closeinnerspacing
punctordspacing     punctopspacing      punctbinspacing   punctrelspacing
punctopenspacing    punctclosespacing   punctpunctspacing punctinnerspacing
innerordspacing     inneropspacing      innerbinspacing   innerrelspacing
inneropenspacing    innerclosespacing   innerpunctspacing innerinnerspacing
\end{typing}

样式参数的值为：

\begin{typing}
display       crampeddisplay
text          crampedtext
script        crampedscript
scriptscript  crampedscriptscript
\end{typing}

值要么是一个数字（表示尺寸或数值），要么是表示 \type {ordordspacing} 及
类似间距参数的 muskip 的黏胶规格节点。

\subsection{{特殊列表头：\type {[get|set}list}}]

\topicindex{lists}

\libindex{lists}
\libindex{setlist}
\libindex{getlist}

虚拟 table \type {tex.lists} 包含跟踪构建页面列表的内部寄存器集合。

\begin{longtable}{l p{5cm}}
\toprule 字段                      & 说明 \\
\midrule
\type{page_ins_head}       & 待处理插入项的循环链表 \\
\type{contrib_head}        & 最近的贡献内容 \\
\type{page_head}           & 当前页面内容 \\
%NC \type{temp_head}           & \\
\type{hold_head}           & 用于为下一页保留的项目 \\
\type{adjust_head}         & 当前 \prm {vadjust} 列表的头 \\
\type{pre_adjust_head}     & 当前 \type {\vadjust pre} 列表的头 \\
%NC \type{align_head}          & \\
\type{page_discards_head}  & 分页断开时被丢弃项目的头 \\
\type{split_discards_head} & vsplit 中被丢弃项目的头 \\
\bottomrule
\end{longtable}

获取和设置函数是 \type {getlist} 和 \type {setlist}。你必须小心设置的内容，
因为 \TEX 可能对列表的构建方式或其所处的状态有期望。

\subsection{语义嵌套层级：\type {getnest}}

\topicindex{nesting}

\libindex{nest}
\libindex{ptr}
%libindex{setnest} % only a message
\libindex{getnest}

虚拟 table \type {nest} 包含当前活动的语义嵌套状态。它有两个主要部分：
一个以零为基的 userdata 数组，表示语义嵌套本身；以及数值 \type {ptr}，
给出最高可用索引。\type {nest[]} 中的数组项和 \type {ptr} 都不能赋值
（因为这会使排版引擎陷入无法修复的混乱），但你可以对数组项内部的各个值
赋值，例如 \type {tex.nest[tex.nest.ptr].prevdepth}。

\type {tex.nest[tex.nest.ptr]} 是当前嵌套状态，\type {nest[0]} 是最外层
（主垂直列表）层级。获取函数是 \type {getnest}。你可以传递一个数字（返回
一个列表）、不传递参数、或 \type {top}（返回最顶层列表），或字符串
\type {ptr}（返回最顶层列表的索引）。

已知字段有：

\begin{longtable}{l l l p{5cm}}
\toprule 键                & 类型    & 模式 & 说明 \\
\midrule
\type{mode}        & number  & 所有   & 这些数字的含义取决于引擎
                                                 有时甚至取决于版本；你可以使用 \typ
                                                 {tex.getmodevalues()} 获取映射：正值
                                                 表示垂直、水平和数学模式，
                                                 负值表示内部和行内变体 \\
\type{modeline}    & number  & 所有   & 进入此模式时的源输入行，
                                                 输出例程内部为负值 \\
\type{head}        & node    & 所有   & 当前列表的头 \\
\type{tail}        & node    & 所有   & 当前列表的尾 \\
\type{prevgraf}    & number  & vmode & 上一段落的行数 \\
\type{prevdepth}   & number  & vmode & 上一段落的深度 \\
\type{spacefactor} & number  & hmode & 当前间距因子 \\
\type{dirs}        & node    & hmode & 供换行算法临时存储使用\\
\type{noad}        & node    & mmode & 用于临时存储待处理的分数分子，
                                                 用于 \prm {over} 等 \\
\type{delimptr}    & node    & mmode & 用于临时存储前一个数学分隔符，
                                                 用于 \prm {middle} \\
\type{mathdir}     & boolean & mmode & 在数学处理期间 \lpr {mathdir} 与
                                                 周围的 \lpr {textdir} 不同时为 true \\
\type{mathstyle}   & number  & mmode & 当前 \lpr {mathstyle} \\
\bottomrule
\end{longtable}

\subsection{打印函数}\label{sec:luaprint}

\topicindex{printing}

\type {tex} table 还包含三个打印函数，它们是从 \LUA 脚本到 \TEX 的主要
接口。这三个函数的参数都存储在一个内存虚拟文件中，该文件作为 \prm {directlua}
展开的结果馈送给 \TEX 扫描器。

一个 \prm {directlua} 命令可返回的文本总量仅受可用系统 \RAM 的限制。不过，
每个单独的打印字符串必须完全适合 \TEX 的输入缓冲区。目前在回调内部使用
这些函数的结果是未定义的。

\subsubsection{\type {print}}

\libindex{print}

\begin{typing}
tex.print(<string> s, ...)
tex.print(<number> n, <string> s, ...)
tex.print(<table> t)
tex.print(<number> n, <table> t)
\end{typing}

每个字符串参数都被 \TEX 视为独立的输入行。如果是 table 参数而不是字符串列表，
则必须是要打印的连续字符串数组（第一个非字符串值将停止打印过程）。

可选参数可用于使用由 \lpr {catcodetable}~\type {n} 定义的类别码规则打印字符串。
如果 \type {n} 为 $-1$，则使用当前活动的类别码规则。如果 \type {n} 为 $-2$，
则生成的类别码是 \prm {the} \prm {toks} 的结果：所有类别码为 12（其他），
除了空格字符的类别码为 10（空格）。否则，如果 \type {n} 不是有效的类别码表，
则被忽略，改为使用当前活动的类别码规则。

在一个 \prm {directlua} 中，最后一个 \type {tex.print} 命令的最后一个字符串
不会附加 \prm {endlinechar}，其他所有字符串都会附加。

\subsubsection{\type {sprint}}

\libindex{sprint}

\begin{typing}
tex.sprint(<string> s, ...)
tex.sprint(<number> n, <string> s, ...)
tex.sprint(<table> t)
tex.sprint(<number> n, <table> t)
\end{typing}

每个字符串参数都被 \TEX 视为一种特殊的输入行，使其适合用作部分行输入机制：

\begin{itemize}
\item \TEX 不会切换到 ``新行 状态，因此前导空格不会被忽略。

\item 不插入 \prm {endlinechar}。

\item 尾部空格不会被移除。注意，这并不妨碍 \TEX 本身在解释行时吞掉空格。
    例如，在

    \begin{typing}
    before\directlua{tex.sprint("\\relax")tex.sprint(" inbetween")}after
    \end{typing}

    中，\type {in between} 之前的空格会因 \prm {relax} 的 ``正常
    扫描而被吞掉。

\end{itemize}

如果是 table 参数而不是字符串列表，则必须是要打印的连续字符串数组（第一个非
字符串值将停止打印过程）。

可选参数设置类别码规则，与 \type {tex.print} 一样。这会影响字符串参数（或
转换为字符串的数字）。

虽然这需要谨慎使用，但你也可以传递令牌或节点 userdata 对象。这些会被注入
流中。令牌最好是有效的令牌，而节点在被注入时必须存在。因此重要的是要认识到
以下几点：

\begin{itemize}\item 当你注入一个令牌时，你需要传递一个有效的令牌 userdata 对象。这个对象在
    不再被引用时会被 \LUA 回收。当它被打印到 \TEX 时，令牌本身会被复制，
    因此不会与 \LUA 垃圾回收产生干扰。你自己管理该对象。由于令牌实际上只是
    数字，在 \TEX 端没有真正的额外开销。

\item 当你注入一个节点时，你需要传递一个有效的节点 userdata 对象。与该对象
    关联的节点在不再被引用时不会被 \LUA 回收。它在 \TEX 端自己的内存空间中
    继续存在。当它被打印到 \TEX 时，节点引用被使用，假设该节点仍然存在。
    不涉及 \LUA 垃圾回收。同样，你自己管理该对象。节点本身在 \TEX 用完后
    被释放。

\end{itemize}

如果你考虑最后一点，可能会意识到当打印的字符串、令牌和节点的混合体被重用时
我们会遇到问题。在 \TEX 内部，序列变为输入缓冲区的链表。因此，\type {"123"}
或 \type {"\foo{123}"} 会被即时读取和解析，而 \typ {<token userdata>} 已经
是令牌化的，实际上现在是一个令牌列表。\typ {<node userdata>} 也被令牌化为
令牌列表，但它有一个对真实节点的引用。通常这没问题。但现在假设你将整个内容
存储在一个宏中：在这种情况下，令牌化节点可以被多次刷新。但是，在第一次这样
的刷新后，节点被使用且其内存被释放。你可以通过使用副本来防止这种情况，这通过
将 \lpr {luacopyinputnodes} 设置为非零值来控制。如果你真的要处理这些底层
问题，这是你必须接受的模糊区域之一。

\subsubsection{\type {tprint}}

\libindex{tprint}

\begin{typing}
tex.tprint({<number> n, <string> s, ...}, {...})
\end{typing}

此函数基本上是重复调用 \type {tex.sprint(<number> n, <string> s, ...)} 的快捷
方式，对每个提供的参数 table 调用一次。

\subsubsection{\type {cprint}}

\libindex{cprint}

此函数接受一个表示要使用的类别码的数字，以及一个字符串 table 或将被推入输入
流的字符串参数列表。

\begin{typing}
tex.cprint( 1," 1: $&{\\foo}") tex.print("\\par") -- a lot of \{ s
tex.cprint( 2," 2: $&{\\foo}") tex.print("\\par") -- matching \} s
tex.cprint( 9," 9: $&{\\foo}") tex.print("\\par") -- all get ignored
tex.cprint(10,"10: $&{\\foo}") tex.print("\\par") -- all become spaces
tex.cprint(11,"11: $&{\\foo}") tex.print("\\par") -- letters
tex.cprint(12,"12: $&{\\foo}") tex.print("\\par") -- other characters
tex.cprint(14,"12: $&{\\foo}") tex.print("\\par") -- comment triggers
\end{typing}

% \subsubsection{\type {write}, \type {twrite}, \type {nwrite}}
\subsubsection{\type {write}}

\libindex{write}
% \libindex{twrite}
% \libindex{nwrite}

\begin{typing}
tex.write(<string> s, ...)
tex.write(<table> t)
\end{typing}

每个字符串参数都被 \TEX 视为一种特殊的输入行，使其适合用作快速转储信息的
方式：

\begin{itemize}\item 该行上的所有类别码要么是 ``空格（用于 '~'）要么是 ``字符（用于其他所有字符）。
\item 不附加 \prm {endlinechar}。
\end{itemize}

如果是 table 参数而不是字符串列表，则必须是要打印的连续字符串数组（第一个非
字符串值将停止打印过程）。

% 函数 \type {twrite} 和 \type {nwrite} 可用于将令牌或节点写回 \TEX，
% 可能与将被令牌化的常规字符串混合。你必须确保传递正确的数据，因为有时
% \TEX 有需要满足的期望。

\subsection{辅助函数}

\subsubsection{\type {round}}

\topicindex {helpers}

\libindex{round}

\begin{typing}
<number> n = tex.round(<number> o)
\end{typing}

对 \LUA 数字 \type {o} 取整，返回一个在有效 \TEX 寄存器值范围内的数字。如果
数字初始超出范围，还会产生 ``number too big 错误。

\subsubsection{\type {scale}}

\libindex{scale}

\begin{typing}
<number> n = tex.scale(<number> o, <number> delta)
<table> n = tex.scale(table o, <number> delta)
\end{typing}

将 \LUA 数字 \type {o} 和 \nod {delta} 相乘，返回一个在有效 \TEX 寄存器值
范围内的取整数字。在 table 版本中，它创建一个 table 的副本，其中所有顶层
数值都按此方式缩放。如果相乘后的数字超出范围，还会产生 ``number too big
错误。

注意：此函数输出的精度取决于你计算机的体系结构和操作系统，因此请谨慎使用！
稍后将添加一个与 \LUATEX 内部 100\% 可移植缩放函数的接口。

\subsubsection{\type {number} 与 \type {romannumeral}}

\libindex{number}
\libindex{romannumeral}

这是原语 \prm {number} 和 \prm {romannumeral} 的配套函数。它们可以这样使用：

\begin{typing}
tex.print(tex.romannumeral(123))
\end{typing}

\subsubsection{\type {fontidentifier} 与 \type {fontname}}

\libindex{fontidentifier}
\libindex{fontname}

第一个仅返回名称，第二个还报告字号。

\begin{typing}
tex.print(tex.fontidentifier(1))
tex.print(tex.fontname(1))
\end{typing}

\subsubsection{\type {sp}}

\libindex{sp}

\begin{typing}
<number> n = tex.sp(<number> o)
<number> n = tex.sp(<string> s)
\end{typing}

将数字 \type {o} 或表示显式尺寸的字符串 \type {s} 转换为缩放点的整数。

对于解析字符串，使用与 \LUATEX 在其类 \TEX 输入语言中扫描尺寸说明符时相同的
扫描和转换规则（包括为错误值生成错误），但以下情况除外：

\begin{itemize}
\item 仅允许显式值，不处理控制序列

\item 无限尺寸单位（\type {fil...}）不允许

\item \type {mu} 单位不产生错误（但可能也没有用处）

\end{itemize}

%\subsubsection{\type {tex.getlinenumber} and \type {tex.setlinenumber}}
%
%\libindex{getlinenumber}
%\libindex{setlinenumber}
%
%你可以操作当前行号：
%
%\begin{typing}
%local n = tex.getlinenumber()
%tex.setlinenumber(n+10)
%\end{typing}
%
%这可以快捷为：
%
%\begin{typing}
%tex.setlinenumber(10,true)
%\end{typing}
%
%当你有一个从文件读取数字并将它们合并在一行中的回调时，这可能很方便（在这种
%情况下，错误消息可能需要引用原始行）。当然，与 \TEX 内部数字处理的干扰是
%可能的。

\subsubsection{\type {error} 与 \type {show_context}}

\topicindex{errors}

\libindex{error}
\libindex{show_context}

\begin{typing}
tex.error(<string> s)
tex.error(<string> s, <table> help)
\end{typing}

这会创建一个错误，类似于 \prm {errhelp} 和 \prm {errmessage} 的组合。在此
错误期间，删除操作被禁用。

\type {help} table 的数组部分必须包含字符串，每行错误帮助一个字符串。

在发生错误时，\type {show_context} 函数将显示我们当前所处的上下文（在展开中）。

\subsubsection{\type {run}、\type {finish}}

\libindex{run}
\libindex{finish}

这两个函数启动解释并强制结束。一次运行通常归结为 \TEX 进入所谓的主循环。
获取一个令牌，根据其当前含义执行某些操作。有时操作立即发生，有时需要更多
扫描。很多时候令牌被推回输入中。这一切意味着 \TEX 扫描器不断推送和弹出输入
状态，但最终在所有操作完成后返回主循环。

\subsubsection{\type {runtoks}}

由于 \TEX 处于展开、处理字体、排版段落、处理 box、构建页面等复杂舞蹈中，
你无法轻易运行嵌套的 \TEX 运行（即嵌套主循环）。不过，有一个选项可以通过
\type {runtoks} 强制进行局部运行。给定令牌列表寄存器的内容在本地展开，
然后我们返回到在 \LUA 端触发此展开的位置。也可以传递一个执行某些工作的
函数。你必须确保最后 \TEX 处于健全状态，而这并不总是容易的。更复杂的机制
会使 \TEX 本身复杂化（可能也会损害性能），因此这个简单的局部展开循环
只能将就。

\begin{typing}
tex.runtoks(<token register>)
tex.runtoks(<lua function>)
\end{typing}

当 \prm {tracingnesting} 参数设置为大于~2 的值时，会报告一些关于局部循环
状态的信息。

当传递令牌寄存器时，此函数有两个可选参数：

\begin{typing}
tex.runtoks(<token register>,force,grouped)
\end{typing}

在例如 \type {\edef} 内部，\type {runtoks} 函数的行为（至少尝试）就像它是
一个 \type {\the}。这防止了不良副作用：通常在这样的定义中令牌保持为令牌，
（例如）字符不会变成节点。使用第二个参数，你可以强制局部主循环，无论如何。
第三个参数添加一层分组。

你可以使用 \type {\endlocalcontrol} 或从 \LUA 端使用 \type {tex.quittoks}
退出局部循环。在这种情况下，你结束一层！当然，最终这可能意味着你到达主层级，
在这种情况下额外的结束将触发冗余警告（而不是中止！）。

\subsubsection{\type {forcehmode}}

\libindex{forcehmode}

一个（可能触发错误的）复杂情况的例子是，\TEX 期望处于某种状态，比如水平
模式，你必须确保在从 \LUA 向 \TEX 反馈内容时它确实处于该状态。通常用户
不会遇到问题，但当你开始写入令牌或节点或进行嵌套运行时，可能会出现需要
运行 \type {forcehmode} 的情况。对此没有现成方案，拦截可能的情况会削弱
\LUATEX 的灵活性。

\subsubsection{\type {hashtokens}}

\libindex{hashtokens}

\topicindex{hash}

\begin{typing}
for i,v in pairs (tex.hashtokens()) do ... end
\end{typing}

返回名称列表。这对调试很有用，但注意这也会报告可能因局部重定义而在此刻无法
访问的控制序列：它严格是哈希表的转储。你可以使用 \type {token.create}
检查属性，例如当创建的 table 中的 \type {command} 键等于 \type {123} 时，
你的 \type {cmdname} 值为 \type {undefined_cs}。

\subsubsection{\type {definefont}}

\topicindex{fonts+defining}

\libindex{definefont}

\begin{typing}
tex.definefont(<string> csname, <number> fontid)
tex.definefont(<boolean> global, <string> csname, <number> fontid)
\end{typing}

将 \type {csname} 与内部字体编号 \type {fontid} 关联。当且仅当指定了
\type {global} 且为 true 时，定义是全局的（不考虑 \type {globaldefs} 的
设置）。

\subsection{处理原语的函数}\label{luaprimitives}

\subsubsection{\type {enableprimitives}}

\libindex{enableprimitives}

\topicindex{initialization}
\topicindex{primitives}

\begin{typing}
tex.enableprimitives(<string> prefix, <table> primitive names)
\end{typing}

此函数接受一个前缀字符串和一个原语名称数组。对于 ``prefix 和
``name 的每个组合，\type {tex.enableprimitives} 首先验证
``name 是否是实际的原语（它必须由下面解释的 \type {tex.extraprimitives}
调用之一返回，或是 \TEX82 的一部分，或是 \prm {directlua}）。如果不是，
\type {tex.enableprimitives} 不做任何操作，跳到下一对。

但如果是，它会通过连接 ``prefix 和 ``name 构造一个 csname
变量，除非 ``prefix 已经是 ``name 的实际前缀。在后一种
情况下，它会丢弃 ``prefix，直接使用 ``name。

然后它会检查构造的 csname 是否存在。如果该 csname 当前未定义（注意：这与
\prm {relax} 不同），它将全局定义该 csname 使其具有以下含义：运行属于原语
``name 的代码。如果由于某种原因 csname 已定义，它不做任何操作，
尝试下一对。

示例：

\begin{typing}
tex.enableprimitives('LuaTeX', {'formatname'})
\end{typing}

将定义 \type {\LuaTeXformatname}，其内在含义与文档中说明的原语
\lpr {formatname} 相同，前提是控制序列 \type {\LuaTeXformatname} 当前未定义。

当 \LUATEX 以 \type {--ini} 运行时，只有 \TEX82 原语和 \prm {directlua}
可用，因此没有任何额外原语{\bf 可用}。

如果你想让所有新功能以其默认名称可用（就像现在这样），你必须在格式生成文件
的开头附近添加

\begin{typing}
\ifx\directlua\undefined \else
    \directlua {tex.enableprimitives('',tex.extraprimitives ())}
\fi
\end{typing}

或者你可以根据需要为不同子集选择不同的前缀。

不过，调用某种形式的 \type {tex.enableprimitives} 非常重要，因为如果不这样做，
你最终会得到一个可以运行 \LUA 代码但做不了多少其他事情的 \TEX82 仿制品。
定义的 csname（当然）保存在格式中，并在运行时可用。

\subsubsection{\type {extraprimitives}}

\libindex{extraprimitives}

\begin{typing}
<table> t = tex.extraprimitives(<string> s, ...)
\end{typing}

此函数返回源自所请求字符串值所给定引擎的原语列表。可能的值及其（当前）
返回值在下表中给出。此外，还定义了一些特殊的原语 ``\tex{ }、
``\tex {/} 和 ``\type {-}。

% [Lua code block removed]

\begin{longtable}{l p{8cm}}
\toprule 名称    & 值 \\
\midrule
tex     & [Lua输出] \\
core    & [Lua输出] \\
etex    & [Lua输出] \\
luatex  & [Lua输出] \\
\bottomrule
\end{longtable}

注意，\type {luatex} 不包含 \type {directlua}，因为它被认为是核心原语，
与所有 \TEX82 原语一起，因此它是从 \type {'core'} 返回的列表的一部分。

运行 \type {tex.extraprimitives} 会给你 \type {-ini} 启动时的完整原语列表。
它完全等同于 \type {tex.extraprimitives("etex","luatex")}。

\subsubsection{\type {primitives}}

\libindex{primitives}

\begin{typing}
<table> t = tex.primitives()
\end{typing}

此函数返回 \LUATEX 已知的所有原语列表。

\subsection{核心功能接口}

\subsubsection{\type {badness}}

\libindex{badness}

\begin{typing}
<number> b = tex.badness(<number> t, <number> s)
\end{typing}

此辅助函数在换行计算中很有用。\type {t} 和 \type {s} 是缩放值；函数返回当
总 \type {t} 应由总和为 \type {s} 的量构成时的不良度。返回的数字是
\mathematics {100(t/s)^3} 的合理近似。

\subsubsection{\type {tex.resetparagraph}}

\topicindex {paragraphs+reset}

\libindex{resetparagraph}

此函数重置 \TEX 在遇到新段落时通常会重置的参数。

\subsubsection{\type {linebreak}}

\topicindex {linebreaks}

\libindex{linebreak}

\begin{typing}
local <node> nodelist, <table> info =
    tex.linebreak(<node> listhead, <table> parameters)
\end{typing}

可识别的参数如下：

\begin{longtable}{l l p{5cm}}
\toprule 名称                        & 类型            & 说明 \\
\midrule
\type{pardir}               & string          & \\
\type{pretolerance}         & number          & \\
\type{tracingparagraphs}    & number          & \\
\type{tolerance}            & number          & \\
\type{looseness}            & number          & \\
\type{hyphenpenalty}        & number          & \\
\type{exhyphenpenalty}      & number          & \\
\type{pdfadjustspacing}     & number          & \\
\type{adjdemerits}          & number          & \\
\type{pdfprotrudechars}     & number          & \\
\type{linepenalty}          & number          & \\
\type{lastlinefit}          & number          & \\
\type{doublehyphendemerits} & number          & \\
\type{finalhyphendemerits}  & number          & \\
\type{hangafter}            & number          & \\
\type{interlinepenalty}     & number or table & 如果是 table，则为类似 \prm {interlinepenalties} 的数组 \\
\type{clubpenalty}          & number or table & 如果是 table，则为类似 \prm {clubpenalties} 的数组 \\
\type{widowpenalty}         & number or table & 如果是 table，则为类似 \prm {widowpenalties} 的数组 \\
\type{brokenpenalty}        & number          & \\
\type{emergencystretch}     & number          & 以缩放点为单位 \\
\type{hangindent}           & number          & 以缩放点为单位 \\
\type{hsize}                & number          & 以缩放点为单位 \\
\type{leftskip}             & glue_spec 节点  & \\
\type{rightskip}            & glue_spec 节点  & \\
\type{parshape}             & table           & \\
\bottomrule
\end{longtable}

注意，\prm {displaywidowpenalties} 没有接口，你必须自己为 \type {widowpenalties}
传递正确的选择。

确保 \type {listhead} 是一个正确的段落列表是你自己的责任：此函数不会向其
添加任何节点。确切地说，如果你想替换核心换行，你可能需要执行以下操作（当你
不在 \cbk {pre_linebreak_filter} 或 \cbk {linebreak_filter} 回调中工作时，
或者当从 listhead 开始的原始列表是在水平模式下生成时）：

\begin{itemize}\item 在开头添加一个 ``缩进 box，可能还有一个 \nod {local_par} 节点
    （仅当你需要时）

\item 将找到的任何最终黏胶替换为无限惩罚（或者如果最后一个节点不是黏胶，
    则添加这样的惩罚）

\item 在该惩罚节点之后为 \prm {parfillskip} 添加一个黏胶节点

\item 确保所有 \type {prev} 指针正确

\end{itemize}

结果是一个节点列表，如果你想将其赋值给 \prm {vbox}，它仍然需要 vpack。
返回的 \type {info} table 包含四个值，全部为数字：

\begin{longtable}{l p{5cm}}
\toprule 名称      & 说明 \\
\midrule
prevdepth & 断开段落中最后一行的深度 \\
prevgraf  & 断开段落中的行数 \\
looseness & 断开段落中的实际 looseness 值 \\
demerits  & 所选方案的总不良度  \\
\bottomrule
\end{longtable}

注意，有几件事你无法通过此函数进行接口：除了通过 \type {pdfadjustspacing}
之外，你无法影响字体扩展，因为其设置在其他地方进行。hbadness 和 hfuzz 等
也是如此。所有这些都在 \type {hpack} 例程中，它通过全局变量获取自己的变量。

\subsubsection{\type {shipout}}

\topicindex {shipout}

\libindex{shipout}

\begin{typing}
tex.shipout(<number> n)
\end{typing}

将 box 编号 \type {n} 发货到输出文件，并清除该 box 寄存器。

\subsubsection{\type {getpagestate}}

\topicindex {pages}

\libindex{getpagestate}

此辅助函数报告当前页面状态：\type {empty}、\type {box_there} 或
\type {inserts_only}，以整数值表示。

\subsubsection{\type {getlocallevel}}

\topicindex {nesting}

\libindex{getlocallevel}

此整数报告局部循环的当前层级。它仅用于调试，（相对状态）数字可能随实现而
变化。

\subsection{随机数生成器}

\libindex{lua_math_random}
\libindex{lua_math_randomseed}
\libindex{init_rand}
\libindex{normal_rand}
\libindex{uniform_rand}
\libindex{uniformdeviate}

由于实际原因，\LUATEX 有自己的随机数生成器。原始的 \LUA 随机函数以
\typ {tex.lua_math_random} 可用。你可以使用新种子通过 \type {init_rand}
初始化（\typ {lua_math_randomseed} 与此等价）。

有三个生成器：\type {normal_rand}（不使用参数）、\type {uniform_rand}
（接受一个在使用前会被取整的数字）和 \type {uniformdeviate}，其行为类似于
原语，期望一个缩放整数，因此

\begin{typing}
tex.print(tex.uniformdeviate(65536)/65536)
\end{typing}

将给出一个介于零和一之间的随机数。

\subsection{与 synctex 相关的函数}\label{synctex}

\topicindex {synctex}

\libindex{set_synctex_mode}      \libindex{get_synctex_mode}
\libindex{set_synctex_no_files}
\libindex{set_synctex_tag}       \libindex{get_synctex_tag}   \libindex{force_synctex_tag}
\libindex{set_synctex_line}      \libindex{get_synctex_line}  \libindex{force_synctex_line}

以下辅助函数仅在你实现自己的 synctex 逻辑时才有意义。请记住，编辑器中使用的
库假定某种逻辑，并且面向 plain 和 \LATEX，因此经过十年，用户期望某种行为。

\begin{longtable}{l p{5cm}}
\toprule 名称                        & 说明 \\
\midrule
\type{set_synctex_mode}     & \type {0} 为默认值，使用正常 synctex
                                    逻辑，\type {1} 使用以下辅助函数设置的值，
                                    而 \type {2} 还为字形节点设置这些值；
                                    \type{3} 设置字形和黏胶，\type {4}
                                    仅设置字形 \\
\type{set_synctex_tag}      & 设置当前标签（文件）值（遵循保存栈） \\
\type{set_synctex_line}     & 设置当前行值（遵循保存栈） \\
\type{set_synctex_no_files} & 禁用 synctex 文件日志记录 \\
\type{get_synctex_mode}     & 返回当前模式（值见上文） \\
\type{get_synctex_tag}      & 获取当前设置的标签（文件）值 \\
\type{get_synctex_line}     & 获取当前设置的行值 \\
\type{force_synctex_tag}    & 覆盖标签（文件）值（\type {0} 重置） \\
\type{force_synctex_line}   & 覆盖行值（\type {0} 重置） \\
\bottomrule
\end{longtable}

最后一个有些特殊。由于 \SYNCTEX 中文件的注册方式，如果我们提供自己的替代方案
且想避免这种开销，我们需要显式禁用该功能。传递值 1 将禁用注册。

\section{\type {texconfig}\label{texconfig}}

\topicindex{libraries+\type{texconfig}}

\topicindex {configuration}

这是一个创建为空的 table。启动 \LUA 脚本可以用一些设置填充此 table，
这些设置在加载和执行启动文件后由可执行文件读取。

\begin{longtable}{l l l p{5cm}}
\toprule 键                    & 类型     & 默认值 & 说明 \\
\midrule
\type{kpse_init}      & boolean  & true &
    \type {false} 完全禁用 \KPATHSEA 初始化，并启用以下数值键-值对的解释。
（仅当你实现了{\it 所有}文件查找回调时才取消设置此项！）
\\
    \type{shell_escape} & string & \type {'f'} & 使用 \type {'y'} 或 \type
    {'t'} 或 \type {'1'} 无条件启用 \type {\write18}，\type {'p'}
    启用 \type {shell_escape_commands} 中列出的命令
\\
    shell_escape_commands & string & & 以逗号分隔的命令名称列表，
    即使 \type {shell_escape} 设置为 \type {'p'} 时也可由 \type {\write18}
    执行。逗号周围{\it 不要}使用空格，任何所需的命令参数用空格分隔，
    对任何需要的参数或路径引用使用 \ASCII 双引号（\type {"}）
\\
\type{string_vacancies} & number   &  75000  & 参见 web2c 文档 \\
\type{pool_free}        & number   &   5000  & 参见 web2c 文档 \\
\type{max_strings}      & number   &  15000  & 参见 web2c 文档 \\
\type{strings_free}     & number   &    100  & 参见 web2c 文档 \\
\type{nest_size}        & number   &     50  & 参见 web2c 文档 \\
\type{max_in_open}      & number   &     15  & 参见 web2c 文档 \\
\type{param_size}       & number   &     60  & 参见 web2c 文档 \\
\type{save_size}        & number   &   4000  & 参见 web2c 文档 \\
\type{stack_size}       & number   &    300  & 参见 web2c 文档 \\
\type{dvi_buf_size}     & number   &  16384  & 参见 web2c 文档 \\
\type{error_line}       & number   &     79  & 参见 web2c 文档 \\
\type{half_error_line}  & number   &     50  & 参见 web2c 文档 \\
\type{max_print_line}   & number   &     79  & 参见 web2c 文档 \\
\type{hash_extra}       & number   &      0  & 参见 web2c 文档 \\
\type{pk_dpi}           & number   &     72  & 参见 web2c 文档 \\
\type{trace_file_names} & boolean  & true
&
    \type {false} 禁用 \TEX 正常的文件打开-关闭反馈（假设回调会处理这些）
\\
\type{trace_extra_newline} & boolean  & false & 在 \type{\tracingmacros} 中添加额外换行 \\
\type{file_line_error} & boolean & false
&
    生成 \type {file:line} 样式的错误消息
\\
\type{halt_on_error} & boolean & false
&
    在遇到第一个错误时中止运行
\\
\type{formatname} & string &
&
    如果命令行上未给出格式名称，则将首先测试此键，而不是直接退出
\\
\type{jobname} & string &
&
    如果命令行上未给出输入文件名，则将首先测试此键，而不是直接放弃
\\
\type{level_chr} & number &
&
    放在被跟踪宏前面的字符（见下一个值）
\\
\type{level_max} & number &
&
    当大于零时，设置 \type
    {\tracingmacros} 时将显示输入嵌套层级；超过此值的层级将被裁剪，
    并在前面显示层级
\\
\type{check_dvi_total_pages} & boolean &
&
     在 \DVI 输出模式下，如果为 true，当页数超过 65535 时中止运行。
     这是默认行为。如果为 false，运行将像在 \TEX 中一样继续。
\\
\type{texlua_img} & boolean &
&
     如果为 true，允许在 \TEXLUA 模式下访问 \type {img} 库。如果
     false（默认），\type {img} 库在 \TEXLUA
     模式下不可用（在 1.22.0 之前的 \LUATEX 版本中是无条件可用的）。
     注意，此设置为{\bf 实验性}，可能随时被移除，恕不另行通知。
\\
\bottomrule
\end{longtable}

注意：与 web2c 参数匹配的数值仅在 \type {kpse_init} 显式设置为 \type {false}
时使用。在所有其他情况下，使用 \type {texmf.cnf} 中的正常值。

你可以启动自己的嵌套层级可视化器，例如：

\begin{typing}
callback.register("input_level_string",function(n)
 if tex.tracingmacros > 0 and tex.count.tracingstacklevels > 0 then
    if tex.tracingmacros > 1 then
        return "! " .. string.rep(">",n) .. " "
    end
  end
  return ""
end)
\end{typing}

或者，与其他引擎同步（未检查）：

\newcount\tracingstacklevels

\begin{typing}
\directlua {
  callback.register("input_level_string", function(n)
    if tex.tracingmacros > 0 then
      local l = tex.count.tracingstacklevels
      if l > 0 then
        return string.rep("~",l) .. string.rep(".",n-l)
      end
    end
    return ""
  end)
}
\end{typing}

\section{\type {texio}}

\topicindex{libraries+\type{texio}}
\topicindex{I/O}

此库负责底层 I/O 接口：写入日志文件和/或控制台。

\subsection{\type {write}}

\libindex{write}

\begin{typing}
texio.write(<string> target, <string> s, ...)
texio.write(<string> s, ...)
\end{typing}

不带 \type {target} 参数时，将所有给定字符串写入 \TEX 此刻写入消息的相同位置。
如果 \prm {batchmode} 生效，则仅写入日志，否则写入日志和终端。可选的
\type {target} 可以是三种可能性之一：\type {term}、\type {log} 或
\type {term and log}。

注意：如果给定多个字符串，且第一个字符串是或可能是上述目标之一，则必须显式
指定 \type {target}，以防止 \LUA 将第一个字符串解释为目标。

\subsection{\type {write_nl}}

\libindex{write_nl}

\begin{typing}
texio.write_nl(<string> target, <string> s, ...)
texio.write_nl(<string> s, ...)
\end{typing}

此函数的行为类似于 \type {texio.write}，但确保给定字符串出现在新行的开头。
如果你只想移到下一行，可以传递一个空字符串。

\subsection{\type {setescape}}

\libindex{setescape}

你可以通过传递零值来禁用控制字符的 \type {\textasciicircum\textasciicircum{}} 转义。

\subsection{\type {closeinput}}

\libindex{closeinput}

此函数应谨慎使用。它的作用类似于 \prm {endinput}，但在 \LUA 端。你可以
使用它（某种程度上）强制跳回 \TEX。通常 \LUA 只会收集打印内容，最后提升
一个输入层级并刷新这些打印内容。此函数可以帮助你停留在当前层级，但你需要
知道自己在做什么（或者更准确地说：\TEX 对输入做了什么）。

\section{\type {token}}

\subsection{扫描器}

\topicindex{libraries+\type{token}}
\topicindex{tokens}

\libindex{scan_keyword}
\libindex{scan_keyword_cs}
\libindex{scan_int}
\libindex{scan_real}
\libindex{scan_float}
\libindex{scan_dimen}
\libindex{scan_glue}
\libindex{scan_toks}
\libindex{scan_code}
\libindex{scan_string}
\libindex{scan_argument}
\libindex{scan_word}
\libindex{scan_csname}
\libindex{scan_list}

令牌库提供了在 \LUA 级别拦截输入并处理它的手段。该库提供了一个基本扫描器
基础设施，可用于编写接受各种参数的宏。此接口刻意保持通用，性能也相当不错。
可以构建额外的解析器，而不会有太多开销。如何从中受益取决于宏包编写者，因为
\LUATEX 背后的主要原则是提供最小的工具集而非解决方案。扫描函数可能是最
引人入胜的。

\begin{longtable}{l l p{5cm}}
\toprule 函数              & 参数           & 结果 \\
\midrule
\type{scan_keyword}   & string             & 如果给定关键字被吞掉则返回 true；与常规
                                                     \TEX 关键字扫描器一样，这是不区分大小写的
                                                     （且基于 \ASCII） \\
\type{scan_keyword_cs} & string             & 如果给定关键字被吞掉则返回 true；此变体
                                                     区分大小写，也适用于 \UTF8 \\
\type{scan_int}       &                    & 返回一个整数 \\
\type{scan_real}      &                    & 从例如 \type {1}、\type {1.1}、\type {.1} 返回一个数字，可带有可选的合并符号 \\
\type{scan_float}     &                    & 从例如 \type {1}、\type {1.1}、\type {.1}、\type {1.1E10}、\type {.1e-10} 返回一个数字，可带有可选的合并符号 \\
\type{scan_dimen}     & infinity、mu-units & 返回一个表示尺寸的数字和/或两个表示填充和阶的数字 \\
\type{scan_glue}      & mu-units           & 返回一个黏胶规格节点 \\
\type{scan_toks}      & definer、expand    & 返回一个令牌 table \\
\type{scan_code}      & bitset             & 如果字符的类别在给定的 bitset（表示类别码）中则返回该字符 \\
\type{scan_string}    &                    & 返回 \type {{}} 之间给出的字符串，作为 \type {\macro} 或作为类别码 11 或 12 的字符序列 \\
\type{scan_argument}  & boolean            & 此函数类似于 \type {scanstring}，但也接受 \type {\cs} \\
\type{scan_word}      &                    & 返回类别码为 11 或 12 的字符序列作为字符串 \\
\type{scan_csname}    &                    & 扫描 \type {\foo} 后返回 \type {foo} \\
\type{scan_list}      &                    & 获取一个 box 规范并返回一个 \type {[h|v]list} 节点 \\
\bottomrule
\end{longtable}

除了扫描令牌的那个之外，扫描器可以被认为是稳定的。\type {scan_code} 函数
接受一个可选数字，\type {keyword} 函数接受一个正常的 \LUA 字符串。
\type {infinity} 布尔值表示我们还允许 \type {fill} 作为尺寸，\type {mu-units}
标志扫描器我们期望数学单位。扫描令牌时，我们可以指示正在定义一个宏，在这种
情况下结果还会提供关于期望什么参数的信息，在结果中这通过一个分隔令牌与含义
分开。\type {expand} 标志决定列表是否会被展开。

\type {scan_argument} 函数展开给定的参数。当扫描带花括号的参数时，可以通过
传递 \type {false} 来禁止展开（默认为 \type {true}）。对于控制序列，传递
\type {false} 将导致一级展开（宏的含义）。

字符串扫描器扫描花括号之间的内容并在过程中展开，或者当看到控制序列时返回
其含义。否则它会扫描类别码为 \type {letter} 或 \type {other} 的字符。因此，
给定以下定义：

\begin{typing}
\def\bar{bar}
\def\foo{foo-\bar}
\end{typing}
bar        : \type{bar}
foo        : \type{foo}
myfirstbar : \type{myfirstbar}

\begin{flushleft}

\end{flushleft}

这是按以下方式输入的：



你看到报告的是 \type {bar} 而不是 \type {myfirstbar} 的原因是，
\type {\bar} 已在之前的段落中使用过。

如果我们现在说：

\begin{typing}
\def\foo{}
\end{typing}

如果我们说

\begin{typing}
\def\foo{\bar}
\end{typing}

当从 \LUA 扫描时，我们根本不处于定义（未定义）宏的模式。在那里我们只得到
真正的原始未定义宏令牌。

\begin{typing}
\directlua{local t = token.get_next() tex.print(t.id.." "..t.tok)}\myfirstbar
\directlua{local t = token.get_next() tex.print(t.id.." "..t.tok)}\mysecondbar
\directlua{local t = token.get_next() tex.print(t.id.." "..t.tok)}\mythirdbar
\end{typing}

\begin{flushleft}

因此，我们确实得到了一个唯一的令牌，因为毕竟我们需要某种可使用和垃圾回收的
\LUA 对象，但它基本上是同一个，表示一个未定义的控制序列。
\end{flushleft}

\section{\type {kpse}}

\topicindex{libraries+\type{kpse}}

此库提供两个独立但几乎相同的 \KPATHSEA 文件搜索功能接口：有一个 ``正常的过程式接口，与 \LUATEX 本身共享其 kpathsea 实例；还有一个完全面向
对象的接口。

\subsection{\type {set_program_name}}

\libindex{set_program_name}
\libindex{default_texmfcnf}
\libindex{new}

库查找变量的方式由 \type {texmf.cmf} 文件驱动，其中当前设置的程序名称用作
过滤器。你可以通过 \type {default_texmfcnf} 检查使用了哪个文件。

在使用搜索库之前，必须先初始化其数据库。有三种可能性，其中两种属于过程式
接口。

首先，当 \LUATEX 用于排版文档时，此初始化自动发生，\KPATHSEA 可执行文件
和程序名称被设置为 \type {luatex}（即，除非用户的启动脚本明确禁止。
详见~\ref{init}）。

其次，在 \TEXLUA 模式下，必须通过 \type {kpse.set_program_name} 函数显式
初始化，它设置 \KPATHSEA 可执行文件（以及可选的程序）名称。

\begin{typing}
kpse.set_program_name(<string> name)
kpse.set_program_name(<string> name, <string> progname)
\end{typing}

第二个参数控制 \type {texmf.cnf} 配置文件中 ``点号值的使用，默认为
第一个参数。

第三，如果你更喜欢面向对象的接口，你必须调用不同的函数。它有相同的参数，
但返回一个 userdata 变量。

\begin{typing}
local kpathsea = kpse.new(<string> name)
local kpathsea = kpse.new(<string> name, <string> progname)
\end{typing}

除了这两个函数之外，接口的调用约定完全相同。根据所选接口，你调用
\type {kpse.find_file} 或 \type {kpathsea:find_file}，具有相同的参数和
返回值。

\subsection{\type {record_input_file}}

\topicindex {files+recording}

\libindex{record_input_file}
\libindex{record_output_file}

这两个函数可用于注册已使用的文件。因为回调可以自己加载文件，你可能需要
这些辅助函数（如果你使用记录功能的话）。

\begin{typing}
kpse.record_input_file(<string> name)
kpse.record_output_file(<string> name)
\end{typing}

\subsection{\type {find_file}}

\topicindex {files+finding}

\libindex {find_file}

库中最常用的函数是 \type {find_file}：

\begin{typing}
<string> f = kpse.find_file(<string> filename)
<string> f = kpse.find_file(<string> filename, <string> ftype)
<string> f = kpse.find_file(<string> filename, <boolean> mustexist)
<string> f = kpse.find_file(<string> filename, <string> ftype, <boolean> mustexist)
<string> f = kpse.find_file(<string> filename, <string> ftype, <number> dpi)
\end{typing}

参数：

\begin{itemize}

\sym{filename}

你要查找的文件名，带或不带扩展名。

\sym{ftype}

映射到 \KPSEWHICH 的 \type {-format} 参数。支持的 \type {ftype}
值与独立 \type {kpsewhich} 程序支持的值相同：\startluacode
    local list = {
        "afm",
        "base",
        "bib",
        "bitmap font",
        "bst",
        "cid maps",
        "clua",
        "cmap files",
        "cnf",
        "cweb",
        "dvips config",
        "enc files",
        "fmt",
        "font feature files",
        "gf",
        "graphic|/|figure",
        "ist",
        "lig files",
        "ls-R",
        "lua",
        "map",
        "mem",
        "MetaPost support",
        "mf",
        "mfpool",
        "mft",
        "misc fonts",
        "mlbib",
        "mlbst",
        "mp",
        "mppool",
        "ocp",
        "ofm",
        "opentype fonts",
        "opl",
        "other binary files",
        "other text files",
        "otp",
        "ovf",
        "ovp",
        "pdftex config",
        "pk",
        "PostScript header",
        "subfont definition files",
        "tex",
        "TeX system documentation",
        "TeX system sources",
        "texmfscripts",
        "texpool",
        "tfm",
        "Troff fonts",
        "truetype fonts",
        "type1 fonts",
        "type42 fonts",
        "vf",
        "web",
        "web2c files",
    }
    table.sort(list)
    context("{\\tttf \letterpercent, t}",list)
\stopluacode

默认类型为 \type {tex}。注意：这与 \KPSEWHICH 不同，后者尝试从提供的
扩展名推断文件类型。

\sym{mustexist}

类似于 \KPSEWHICH 的 \type {-must-exist}，默认为 \type {false}。如果你
指定 \type {true}（或非零整数），则 \KPSE 库将同时搜索磁盘和 \type {ls-R}
数据库。

\sym{dpi}

这用于格式 \type {pk}、\type {gf} 和 \type {bitmap font} 的尺寸参数。
\end{itemize}

如果指定了 \type {--output-directory} 且值为相对路径名，则首先在此处搜索
文件，如果失败，则在标准树中搜索。

\subsection{\type {lookup}}

\libindex{lookup}

还有一种更强大（但更慢）的通用文件查找方法。它为每个找到的文件返回一个
字符串。

\begin{typing}
<string> f, ... = kpse.lookup(<string> filename, <table> options)
\end{typing}

选项与 \type {kpsewhich} 的命令行参数对应：

\begin{longtable}{l l p{5cm}}
\toprule 键              & 类型     & 说明 \\
\midrule
\type{debug}     & number   & 为此查找设置调试标志\\
\type{format}    & string   & 使用特定文件类型（见上面的列表）\\
\type{dpi}       & number   & 为此查找使用此分辨率；默认 600\\
\type{path}      & string   & 在给定路径中搜索\\
\type{all}       & boolean  & 输出所有匹配项，而不仅仅是第一个\\
\type{mustexist} & boolean  & 必要时同时搜索磁盘和 ls-R\\
\type{mktexpk}   & boolean  & 为此查找禁用/启用 mktexpk 生成\\
\type{mktextex}  & boolean  & 为此查找禁用/启用 mktextex 生成\\
\type{mktexmf}   & boolean  & 为此查找禁用/启用 mktexmf 生成\\
\type{mktextfm}  & boolean  & 为此查找禁用/启用 mktextfm 生成\\
\type{subdir}    & string
                         或 table & 仅输出目录部分以给定字符串结尾的匹配项 \\
\bottomrule
\end{longtable}

如果指定了 \type {--output-directory} 且值为相对路径名，则首先在此处搜索
文件，然后在标准树中搜索。

\subsection{\type {init_prog}}

\topicindex {initialization+bitmaps}

\libindex{init_prog}

为需要生成位图字体的程序提供额外初始化。

\begin{typing}
kpse.init_prog(<string> prefix, <number> base_dpi, <string> mfmode)
kpse.init_prog(<string> prefix, <number> base_dpi, <string> mfmode, <string> fallback)
\end{typing}

\subsection{\type {readable_file}}

\libindex{readable_file}

测试（绝对）文件名是否为可读文件。

\begin{typing}
<string> f = kpse.readable_file(<string> name)
\end{typing}

返回值是你应该使用的实际绝对文件名，因为磁盘名称并不总是与请求的名称相同，
这是由于别名和例如 \MSDOS 下的系统特定处理。如果文件不存在或不可读则返回
\type {nil}。

\subsection{\type {expand_path}}

\libindex{expand_path}

类似于 kpsewhich 的 \type {-expand-path}：

\begin{typing}
<string> r = kpse.expand_path(<string> s)
\end{typing}

\subsection{\type {expand_var}}

\libindex{expand_var}

类似于 kpsewhich 的 \type {-expand-var}：

\begin{typing}
<string> r = kpse.expand_var(<string> s)
\end{typing}

\subsection{\type {expand_braces}}

\libindex{expand_braces}

类似于 kpsewhich 的 \type {-expand-braces}：

\begin{typing}
<string> r = kpse.expand_braces(<string> s)
\end{typing}

\subsection{\type {in_name_ok}}

\libindex{in_name_ok}

如果 \type{fname} 可打开读取则返回 true，否则返回 false 并向标准错误写入
一条消息。

\begin{typing}
<boolean> r = kpse.in_name_ok(<string> fname)
\end{typing}

\subsection{\type {in_name_ok_silent_extended}}

\libindex{in_name_ok_silent_extended}

\begin{typing}
<boolean> r = kpse.in_name_ok_silent_extended(<string> fname)
\end{typing}

如果 \type{fname} 可打开读取则返回 true；
还会检查 \type{TEXMFVAR} 和 \type{TEXMFSYSVAR} 的值是否为绝对文件名。
否则返回 false，但不向标准错误写入消息。

\subsection{\type {out_name_ok}}

\libindex{out_name_ok}

如果 \type{fname} 可打开写入则返回 true。

\begin{typing}
<boolean> r = kpse.out_name_ok(<string> fname)
\end{typing}

\subsection{\type {out_name_ok_silent_extended}}

\libindex{out_name_ok_silent_extended}

\begin{typing}
<boolean> r = kpse.out_name_ok_silent_extended(<string> fname)
\end{typing}

如果 \type{fname} 可打开写入则返回 true；
还会检查 \type{TEXMFVAR} 和 \type{TEXMFSYSVAR} 的值是否为绝对文件名。
否则返回 false，但不向标准错误写入消息。

\subsection{\type {show_path}}

\libindex{show_path}

类似于 kpsewhich 的 \type {-show-path}：

\begin{typing}
<string> r = kpse.show_path(<string> ftype)
\end{typing}

\subsection{\type {var_value}}

\libindex{var_value}

类似于 kpsewhich 的 \type {-var-value}：

\begin{typing}
<string> r = kpse.var_value(<string> s)
\end{typing}

\subsection{\type {version}}

\libindex{version}

返回 kpathsea 版本字符串。

\begin{typing}
<string> r = kpse.version()
\end{typing}

\subsection{\type {check_permission}}

\libindex{check\_permission}

检查 \type{filename} 是否可执行。
如果可执行则返回 \type{1}、\type{filename} 或安全的 cmd 替代方案，
否则返回 \type{0} 和错误消息。

\begin{typing}
<string> res, cmd = kpse.check_permission(<string> filename)
\end{typing}
